\documentclass[12pt,a4paper]{article}
\usepackage[utf8]{inputenc}
\usepackage[T1]{fontenc}
\usepackage{amsmath}
\usepackage{amssymb}
\usepackage{graphicx}
\usepackage[spanish]{babel}
\usepackage[a4paper, margin=2.5cm]{geometry}
\usepackage[european]{circuitikz}
\usetikzlibrary{babel} % Para compatibilidad con babel-spanish
\usepackage{comment} % Para poder usar entornos de comentario

\title{Análisis Nodal del Circuito}
\author{Gemini Code Assist and Prof. CERO}
\date{\today}

\begin{document}

\maketitle

\section{Descripción del Circuito}

El circuito a analizar consta de tres ramas en paralelo conectadas entre un nodo superior, que denominaremos $V_x$, y tierra (GND). El objetivo es determinar la fórmula para calcular la tensión $V_x$.

\begin{center}
    \begin{circuitikz}[scale=1.2, transform shape]
        \draw (0,0) node[ground] (gnd) {};
        \draw (0,3) node[circ, label=above:$V_x$] (vx) {};
        
        % Rama Izquierda (R1 e Is)
        \draw (vx) to[R, l=$R_1$] (-3,3) to[I, i<=$I_s$] (-3,0) -- (gnd);
        
        % Rama Central (R3)
        \draw (vx) to[R, l=$R_3$] (0,0);
        
        % Rama Derecha (R2 y E3)
        \draw (vx) to[R, l=$R_2$] (3,3) to[V, v^<=$E_3$] (3,0) -- (gnd);
    \end{circuitikz}
\end{center}

Los valores de los componentes son:
\begin{itemize}
    \item Resistencia $R_1 = 50 \, \Omega$
    \item Resistencia $R_2 = 50 \, \Omega$
    \item Resistencia $R_3 = 200 \, \Omega$
    \item Fuente de corriente $I_s = 50 \, \text{mA} = 0.05 \, \text{A}$
    \item Fuente de tensión $E_3 = 50 \, \text{V}$
\end{itemize}

\section{Derivación de la Fórmula para $V_x$}

Para encontrar la tensión en el nodo $V_x$, aplicamos la Ley de Corrientes de Kirchhoff (LCK), estableciendo que la suma de las corrientes que salen del nodo es igual a cero.

\begin{equation}
    I_{R3} + I_{R2} - I_s = 0
\end{equation}

Donde la corriente que sale del nodo $V_x$ por cada rama es:
\begin{itemize}
    \item \textbf{Rama Izquierda:} La corriente que llega al nodo es $I_s$, por lo que la corriente que sale es $-I_s$. La resistencia $R_1$ no influye en la tensión del nodo $V_x$ al estar en serie con una fuente de corriente ideal.
    \item \textbf{Rama Central:} La corriente que sale por $R_3$ es $I_{R3} = \frac{V_x}{R_3}$.
    \item \textbf{Rama Derecha:} La corriente que sale por $R_2$ es $I_{R2} = \frac{V_x - E_3}{R_2}$. El potencial en el terminal superior de la fuente $E_3$ es $+E_3$ con respecto a tierra.
\end{itemize}

Sustituyendo estas expresiones en la ecuación de la LCK:
\begin{equation}
    \frac{V_x}{R_3} + \frac{V_x - E_3}{R_2} - I_s = 0
\end{equation}

Agrupamos los términos que contienen $V_x$:
\begin{equation}
    V_x \left( \frac{1}{R_2} + \frac{1}{R_3} \right) - \frac{E_3}{R_2} - I_s = 0
\end{equation}

Finalmente, despejamos $V_x$ para obtener la fórmula general:
\begin{equation}
    V_x = \frac{I_s + \frac{E_3}{R_2}}{\frac{1}{R_2} + \frac{1}{R_3}}
\end{equation}

\section{Cálculo del Valor de $V_x$}

Sustituyendo los valores numéricos del circuito en la fórmula obtenida:
\begin{align*}
    V_x &= \frac{0.05 \, \text{A} + \frac{50 \, \text{V}}{50 \, \Omega}}{\frac{1}{50 \, \Omega} + \frac{1}{200 \, \Omega}} \\
    V_x &= \frac{0.05 + 1}{0.02 + 0.005} \\
    V_x &= \frac{1.05}{0.025} \\
    V_x &= 42 \, \text{V}
\end{align*}

La tensión en el punto de unión de las tres resistencias es \textbf{42 V}, el
cual coincide con el reultado que arroja el simulador (\textbf{DroidTesla}).

\end{document}